\documentclass[11pt, a4paper]{report}

% ========== Common Packages ==========
\usepackage{fontspec}
\usepackage{kotex}
\usepackage{amsmath, amssymb, amsthm}
\usepackage{mathtools}
\usepackage{bm}
\usepackage{graphicx}
\usepackage{booktabs}
\usepackage{array}
\usepackage{tabularx}
\usepackage{longtable}
\usepackage{hyperref}
\usepackage{cleveref}
\usepackage{geometry}
\usepackage{fancyhdr}
\usepackage{titlesec}
\usepackage{enumitem}
\usepackage{xcolor}
\usepackage{tcolorbox}
\usepackage{algorithm}
\usepackage{algorithmic}
\usepackage{tikz}
\usetikzlibrary{arrows.meta, positioning, calc, decorations.pathreplacing, shapes.geometric, fit, patterns, angles, quotes, trees}
\usepackage{pgfplots}
\pgfplotsset{compat=1.18}
\usepackage{caption}
\usepackage{subcaption}
\usepackage{float}

\geometry{margin=2.5cm, top=3cm, bottom=3cm}

% ========== Colors ==========
% Common
\definecolor{defblue}{RGB}{30, 80, 150}
\definecolor{thmgreen}{RGB}{20, 110, 60}
\definecolor{noteorange}{RGB}{200, 100, 0}
\definecolor{lightblue}{RGB}{230, 240, 255}
\definecolor{lightgreen}{RGB}{230, 255, 235}
\definecolor{lightorange}{RGB}{255, 245, 230}
\definecolor{lightgray}{RGB}{245, 245, 245}
% Extended
\definecolor{lightred}{RGB}{255, 235, 235}
\definecolor{darkred}{RGB}{180, 30, 30}
\definecolor{biogreen}{RGB}{10, 130, 80}
\definecolor{cvpurple}{RGB}{100, 40, 150}
\definecolor{injuryred}{RGB}{200, 40, 40}
\definecolor{codegray}{RGB}{240, 240, 245}
\definecolor{pipeblue}{RGB}{25, 100, 180}
\definecolor{constraintred}{RGB}{200, 50, 50}
\definecolor{termblack}{RGB}{30, 30, 30}
\definecolor{goldcolor}{RGB}{180, 140, 20}

% ========== Theorem-like environments ==========
\theoremstyle{definition}
\newtheorem{definition}{Definition}[chapter]
\newtheorem{example}{Example}[chapter]

\theoremstyle{plain}
\newtheorem{theorem}{Theorem}[chapter]
\newtheorem{proposition}{Proposition}[chapter]

\theoremstyle{remark}
\newtheorem{remark}{Remark}[chapter]

% ========== tcolorbox environments ==========
% --- Common (all parts) ---
\newtcolorbox{keyidea}[1][]{
  colback=lightblue,
  colframe=defblue,
  fonttitle=\bfseries,
  title={Key Idea},
  #1
}

\newtcolorbox{importantnote}[1][]{
  colback=lightorange,
  colframe=noteorange,
  fonttitle=\bfseries,
  title={Important},
  #1
}

\newtcolorbox{summary}[1][]{
  colback=lightgreen,
  colframe=thmgreen,
  fonttitle=\bfseries,
  title={Summary},
  #1
}

% --- Warning/Error ---
\newtcolorbox{warningbox}[1][]{
  colback=lightred,
  colframe=darkred,
  fonttitle=\bfseries,
  title={Warning},
  #1
}

% --- Domain: Biomechanics & Clinical ---
\newtcolorbox{clinicalnote}[1][]{
  colback=lightred,
  colframe=darkred,
  fonttitle=\bfseries,
  title={Clinical Note},
  #1
}

\newtcolorbox{biomechbox}[1][]{
  colback=lightgreen,
  colframe=biogreen,
  fonttitle=\bfseries,
  title={Biomechanics},
  #1
}

\newtcolorbox{injurybox}[1][]{
  colback={injuryred!8},
  colframe=injuryred,
  fonttitle=\bfseries,
  title={Injury Risk},
  #1
}

% --- Domain: Computer Vision ---
\newtcolorbox{cvbox}[1][]{
  colback={cvpurple!5},
  colframe=cvpurple,
  fonttitle=\bfseries,
  title={Computer Vision},
  #1
}

% --- Pipeline & Setup ---
\newtcolorbox{setupbox}[1][]{
  colback=lightgray,
  colframe=gray!60,
  fonttitle=\bfseries,
  title={Setup},
  #1
}

\newtcolorbox{pipelinebox}[1][]{
  colback=pipeblue!5,
  colframe=pipeblue,
  fonttitle=\bfseries,
  title={Pipeline Step},
  #1
}

\newtcolorbox{codebox}[1][]{
  colback=codegray,
  colframe=gray!70,
  fonttitle=\bfseries\ttfamily,
  title={Code},
  #1
}

\newtcolorbox{termbox}[1][]{
  colback=termblack!5,
  colframe=termblack!60,
  fonttitle=\bfseries\ttfamily,
  title={Terminal},
  #1
}

% --- Research ---
\newtcolorbox{hypothesis}[1][]{
  colback=goldcolor!8,
  colframe=goldcolor,
  fonttitle=\bfseries,
  title={Hypothesis},
  #1
}

\newtcolorbox{resultbox}[1][]{
  colback=thmgreen!8,
  colframe=thmgreen,
  fonttitle=\bfseries,
  title={Expected Result},
  #1
}

% --- Constraints & Solutions ---
\newtcolorbox{constraintbox}[1][]{
  colback=constraintred!8,
  colframe=constraintred,
  fonttitle=\bfseries,
  title={Constraint},
  #1
}

\newtcolorbox{solutionbox}[1][]{
  colback=thmgreen!8,
  colframe=thmgreen,
  fonttitle=\bfseries,
  title={Solution},
  #1
}

\newtcolorbox{databox}[1][]{
  colback=pipeblue!5,
  colframe=pipeblue,
  fonttitle=\bfseries,
  title={Data Spec},
  #1
}

% --- Progress/Status ---
\newtcolorbox{successbox}[1][]{
  colback=lightgreen,
  colframe=thmgreen,
  fonttitle=\bfseries,
  title={Success},
  #1
}

\newtcolorbox{nextbox}[1][]{
  colback=lightorange,
  colframe=noteorange,
  fonttitle=\bfseries,
  title={Next Step},
  #1
}

% --- Fitting guide ---
\newtcolorbox{failbox}[1][]{
  colback=lightred,
  colframe=darkred,
  fonttitle=\bfseries,
  title={Failure Pattern},
  #1
}

\newtcolorbox{fixbox}[1][]{
  colback=lightgreen,
  colframe=thmgreen,
  fonttitle=\bfseries,
  title={Fix},
  #1
}

\newtcolorbox{lessonbox}[1][]{
  colback=lightorange,
  colframe=noteorange,
  fonttitle=\bfseries,
  title={Lesson Learned},
  #1
}

% --- Specification ---
\newtcolorbox{specbox}[1][]{
  colback=cvpurple!5,
  colframe=cvpurple,
  fonttitle=\bfseries,
  title={Specification},
  #1
}

\newtcolorbox{checklistbox}[1][]{
  colback=lightgreen,
  colframe=biogreen,
  fonttitle=\bfseries,
  title={Checklist},
  #1
}

% ========== Custom math commands ==========
% Vectors (lowercase)
\newcommand{\vx}{\mathbf{x}}
\newcommand{\vd}{\mathbf{d}}
\newcommand{\vo}{\mathbf{o}}
\newcommand{\vc}{\mathbf{c}}
\newcommand{\vr}{\mathbf{r}}
\newcommand{\vp}{\mathbf{p}}
\newcommand{\vv}{\mathbf{v}}
\newcommand{\vn}{\mathbf{n}}
\newcommand{\vu}{\mathbf{u}}
\newcommand{\vt}{\mathbf{t}}
\newcommand{\ve}{\mathbf{e}}
\newcommand{\vq}{\mathbf{q}}
% Vectors (uppercase)
\newcommand{\vF}{\mathbf{F}}
\newcommand{\vM}{\mathbf{M}}
\newcommand{\va}{\mathbf{a}}
\newcommand{\vg}{\mathbf{g}}
\newcommand{\vX}{\mathbf{X}}
% Bold Greek
\newcommand{\vmu}{\bm{\mu}}
\newcommand{\vbeta}{\bm{\beta}}
\newcommand{\vtheta}{\bm{\theta}}
\newcommand{\vpsi}{\bm{\psi}}
% Matrices
\newcommand{\mSigma}{\bm{\Sigma}}
\newcommand{\mR}{\mathbf{R}}
\newcommand{\mS}{\mathbf{S}}
\newcommand{\mW}{\mathbf{W}}
\newcommand{\mJ}{\mathbf{J}}
\newcommand{\mG}{\mathbf{G}}
\newcommand{\mT}{\mathbf{T}}
\newcommand{\mI}{\mathbf{I}}
\newcommand{\mB}{\mathbf{B}}
\newcommand{\mK}{\mathbf{K}}
\newcommand{\mP}{\mathbf{P}}
\newcommand{\mF}{\mathbf{F}}
\newcommand{\mE}{\mathbf{E}}
\newcommand{\mA}{\mathbf{A}}
\newcommand{\mH}{\mathbf{H}}
% Sets and groups
\newcommand{\RR}{\mathbb{R}}
\newcommand{\SO}{\mathrm{SO}}
% Units
\newcommand{\degps}{\,\text{deg/s}}
\newcommand{\Nm}{\ensuremath{\,\text{N}\cdot\text{m}}}
\newcommand{\BW}{\,\text{BW}}

% ========== Header/Footer ==========
\pagestyle{fancy}
\fancyhf{}
\fancyhead[L]{\leftmark}
\fancyhead[R]{\thepage}
\renewcommand{\headrulewidth}{0.4pt}

% ========== Title formatting ==========
\titleformat{\chapter}[display]
  {\normalfont\huge\bfseries\color{defblue}}
  {\chaptertitlename\ \thechapter}{20pt}{\Huge}

\titleformat{\section}
  {\normalfont\Large\bfseries\color{defblue!80!black}}
  {\thesection}{1em}{}

\titleformat{\subsection}
  {\normalfont\large\bfseries\color{defblue!60!black}}
  {\thesubsection}{1em}{}


\begin{document}

% ==================== TITLE PAGE ====================
\begin{titlepage}
\centering
\vspace*{2cm}
{\Huge\bfseries\color{defblue} 7-View 마커리스 투구 분석\\[0.3cm] 전체 실행 계획서}

\vspace{1cm}
{\LARGE\color{defblue!70!black} 소계획 $\to$ 중계획 $\to$ 대계획 $\to$ 전체 계획}

\vspace{2cm}
{\large 파이프라인 요소별 완전 분석 + GART 포토메트릭 정제}

\vspace{1cm}
{\large 2026년 4월 14일}

\vspace{3cm}
\begin{tcolorbox}[colback=lightblue, colframe=defblue, width=0.85\textwidth]
\textbf{목표}: 7개 카메라 전체 활용, SMPL reproj 오차 153px $\to$ $<$50px,\\
GART 30,000스텝 학습, 포토메트릭 포즈 정제로 모션 품질 최대화
\end{tcolorbox}
\end{titlepage}

\tableofcontents
\newpage

% ====================================================================
% CHAPTER 1: 요소별 소계획 (11개)
% ====================================================================
\chapter{요소별 소계획}

각 파이프라인 요소를 독립적으로 분석하고, 해당 요소의 개선 소계획을 수립합니다.

% ---------------------------------------------------------------
\section{소계획 1: 카메라 캘리브레이션}
\label{sec:sp1}

\subsection{현재 상태}
\begin{center}
\renewcommand{\arraystretch}{1.2}
\begin{tabular}{lcccc}
\toprule
\textbf{카메라} & \textbf{해상도} & \textbf{방향} & \textbf{FOV (H)} & \textbf{COLMAP 등록} \\
\midrule
cam1 & 1080$\times$1920 & 세로 & 72.7° & \checkmark \\
cam2 & 1920$\times$1080 & 가로 & 104.4° & \checkmark \\
cam3 & 1920$\times$1080 & 가로 & 104.7° & \checkmark \\
cam4 & 1920$\times$1080 & 가로 & 122.1° & \checkmark \\
cam5 & 1920$\times$1080 & 가로 & 117.2° & \checkmark \\
cam6 & 1920$\times$1080 & 가로 & 70.2° & \checkmark \\
cam7 & 1920$\times$1080 & 가로 & 99.6° & \checkmark \\
\bottomrule
\end{tabular}
\end{center}

\begin{itemize}
\item COLMAP Round 2: 336/336 등록 (100\%), 중간 reproj 오차 0.75px
\item SIMPLE\_RADIAL 모델 (f, cx, cy, k1)
\item VGGT 캘리브레이션도 별도 완료 (K, R, t at full resolution)
\end{itemize}

\subsection{문제점}
\begin{enumerate}
\item cam1 세로 방향 $\to$ U-HMR 입력 불가 (1080$\times$1920)
\item VGGT vs COLMAP 좌표계 불일치 가능성
\item cam4 FOV 122° --- 강한 왜곡 가능
\end{enumerate}

\subsection{소계획}
\begin{pipelinebox}[title={SP-1: 캘리브레이션 7-view 통합}]
\begin{enumerate}
\item cam1 이미지 90° CW 회전 $\to$ 1920$\times$1080 가로 변환
\item cam1 K행렬 변환: $f_x \leftrightarrow f_y$, $c_x \leftrightarrow c_y$ 교환 + 좌표축 반영
\item COLMAP Round 2 파라미터를 기준으로 7개 카메라 통합 K, R, t 추출
\item ZJU-MoCap \texttt{annots.npy} 형식으로 변환
\item 왜곡 보정 (cam4 k1 적용 여부 결정)
\end{enumerate}
\textbf{산출물}: \texttt{annots\_7view.npy} (7개 카메라 K, R, t, D)
\end{pipelinebox}


% ---------------------------------------------------------------
\section{소계획 2: 2D 포즈 추정 (OpenPose)}
\label{sec:sp2}

\subsection{현재 상태}
\begin{itemize}
\item 7개 카메라 $\times$ 1000 프레임 OpenPose Body25 추출 완료
\item 4단계 스무딩 (v2): 83--87\% 지터 감소
\item 신뢰도 임계값: 0.3
\end{itemize}

\subsection{문제점}
\begin{enumerate}
\item 가속기(30ms)에서 모션 블러 $\to$ 손목/팔꿈치 신뢰도 급감
\item cam1 세로 이미지에서 OpenPose 미실행 가능성
\item ViTPose cam3 91\% 오인식 문제 (OpenPose로 우회 완료)
\end{enumerate}

\subsection{소계획}
\begin{pipelinebox}[title={SP-2: 7-view 2D 키포인트 완성}]
\begin{enumerate}
\item cam1 회전 이미지에 OpenPose 실행 (1000 프레임)
\item 기존 cam2--cam7 스무딩 결과 재활용 (변경 불필요)
\item cam1 결과에 동일 4단계 스무딩 적용
\item 7개 카메라 키포인트를 통합 텐서로 구성: $(1000, 7, 25, 3)$
\end{enumerate}
\textbf{산출물}: 7-view 스무딩된 OpenPose 키포인트
\end{pipelinebox}


% ---------------------------------------------------------------
\section{소계획 3: U-HMR body\_pose 초기값}
\label{sec:sp3}

\subsection{현재 상태}
\begin{itemize}
\item U-HMR Mv\_Fusion: 4-view 고정 아키텍처 (cam3, 5, 6, 7)
\item 999 프레임 완료, A100에서 18분
\item body\_pose: $(23, 3)$ axis-angle, betas: $(10,)$
\item 다시점 융합으로 body\_pose가 4개 뷰에서 동일 (검증됨)
\end{itemize}

\subsection{문제점}
\begin{enumerate}
\item 4-view 제약은 U-HMR 아키텍처 한계 --- 변경 불가
\item body\_pose에 5--15°/관절 내재 오차 존재
\item 프레임별 독립 추론 $\to$ 시간적 비일관성
\end{enumerate}

\subsection{소계획}
\begin{pipelinebox}[title={SP-3: U-HMR 초기값 활용 전략}]
\begin{enumerate}
\item U-HMR 4-view 결과를 \textbf{초기값으로만} 사용 (고정하지 않음)
\item body\_pose 오차는 Stage 2 (7-view reproj)에서 보정
\item betas는 전 프레임 평균으로 고정 (체형은 불변)
\item 누락된 frame 000 보간 (frame 001 복사)
\end{enumerate}
\textbf{산출물}: 1000 프레임 body\_pose 초기값 (변경 없음, 기존 결과 재활용)
\end{pipelinebox}


% ---------------------------------------------------------------
\section{소계획 4: 7-View Reproj Fitting (핵심 개선)}
\label{sec:sp4}

\subsection{현재 상태}
\begin{center}
\renewcommand{\arraystretch}{1.2}
\begin{tabular}{lcc}
\toprule
\textbf{설정} & \textbf{현재 (4-view)} & \textbf{목표 (7-view)} \\
\midrule
카메라 수 & 4 (cam3,5,6,7) & \textbf{7 (cam1--7)} \\
body\_pose 정규화 $\lambda$ & 0.01 & \textbf{0.001} \\
Stage 1 (RT) 스텝/lr & 80 / 0.05 & \textbf{200 / 0.03} \\
Stage 2 (pose) 스텝/lr & 40 / 0.008 & \textbf{300 / 0.01} \\
Stage 3 (betas) & 없음 & \textbf{50 / 0.005 (선택)} \\
Sequential 초기화 & Rh, Th만 & \textbf{Rh, Th, body\_pose 모두} \\
관절각 클램핑 & 없음 & \textbf{팔꿈치 0--150°, 무릎 0--160°} \\
temporal smoothing & 없음 & \textbf{후처리 Savitzky-Golay} \\
\midrule
예상 reproj 오차 & 153.5px & \textbf{$<$50px} \\
\bottomrule
\end{tabular}
\end{center}

\subsection{문제점 (근본 원인 5가지)}
\begin{enumerate}
\item $\lambda=0.01$ body\_pose 정규화 $\to$ 오차 바닥 113px (수정 차단)
\item Stage 2에서 100스텝, lr=0.005 $\to$ 69 파라미터에 부족
\item 4-cam만 사용 $\to$ 3개 카메라 정보 버림 (깊이 제약 약화)
\item Sequential이라고 했지만 Rh, Th만 전달 (body\_pose는 매 프레임 리셋)
\item 관절각 제한 없음 $\to$ 비물리적 해 허용
\end{enumerate}

\subsection{소계획}
\begin{pipelinebox}[title={SP-4: 7-View Sequential Reproj Fitting 스크립트}]
\begin{enumerate}
\item \textbf{7개 카메라 로드}: COLMAP Rd2 K, R, t (cam1 회전 보정 포함)
\item \textbf{정규화 완화}: $\lambda = 0.01 \to 0.001$ (body\_pose 수정 허용)
\item \textbf{Stage 1 강화}: Rh(3) + Th(3), 200스텝, lr=0.03
\item \textbf{Stage 2 강화}: + body\_pose(69), 300스텝, lr=0.01, $\lambda=0.001$
\item \textbf{진짜 Sequential}: frame[i-1]의 (Rh, Th, body\_pose) $\to$ frame[i] 초기값
\item \textbf{관절각 클램핑}: 매 스텝 후 팔꿈치/무릎/어깨 ROM 적용
\item \textbf{GMoF $\sigma$ 적응}: $\sigma = 100 \to 50$ (이상치 더 강하게 억제)
\item \textbf{후처리}: Savitzky-Golay (window=15, poly=3) on all params
\item \textbf{검증}: 프레임별 7-cam reproj 오차 계산 + 히스토그램 저장
\end{enumerate}
\textbf{산출물}: 1000 프레임 정제된 SMPL 파라미터 (Rh, Th, body\_pose, betas)
\end{pipelinebox}


% ---------------------------------------------------------------
\section{소계획 5: 이미지 데이터 준비}
\label{sec:sp5}

\subsection{현재 상태}
\begin{itemize}
\item 7개 카메라 $\times$ 1000 프레임 = 7000 이미지 존재
\item cam1은 세로 (1080$\times$1920), 나머지 가로 (1920$\times$1080)
\item SAM3 마스크 7000장 완료 (투구자 분리)
\end{itemize}

\subsection{소계획}
\begin{pipelinebox}[title={SP-5: ZJU-MoCap 형식 이미지 구조화}]
\begin{enumerate}
\item 디렉토리 구조: \texttt{images/\{0,1,2,3,4,5,6\}/000000.jpg $\sim$ 000999.jpg}
\item cam1 이미지 90° 회전 후 저장 (가로 1920$\times$1080)
\item cam2--7 이미지 심볼릭 링크 또는 복사
\item 배경 마스크 (선택): SAM3 마스크를 GART 학습에 활용 가능
\item 총 용량 확인: 7000 $\times$ $\sim$6MB = $\sim$42GB
\end{enumerate}
\textbf{산출물}: ZJU-MoCap 표준 이미지 디렉토리
\end{pipelinebox}


% ---------------------------------------------------------------
\section{소계획 6: SMPL 파라미터 변환}
\label{sec:sp6}

\subsection{현재 상태}
\begin{itemize}
\item SP-4 출력: per-frame JSON (EasyMocap 형식) + NPZ 아카이브
\item GART 입력: ZJU-MoCap \texttt{smpl\_params/\{000..999\}.npy}
\end{itemize}

\subsection{소계획}
\begin{pipelinebox}[title={SP-6: SMPL 파라미터 $\to$ GART 입력 형식}]
\begin{enumerate}
\item SP-4 출력 (Rh, Th, body\_pose, betas) $\to$ per-frame .npy 변환
\item 각 .npy: \texttt{global\_orient}(3), \texttt{body\_pose}(69), \texttt{transl}(3), \texttt{betas}(10)
\item 좌표계 확인: SMPL Y-up $\leftrightarrow$ COLMAP Y-down 변환 적용
\item \texttt{smpl\_params/} 디렉토리에 1000 파일 저장
\item 무결성 검증: NaN 체크, body\_pose 범위 확인
\end{enumerate}
\textbf{산출물}: \texttt{smpl\_params/000.npy $\sim$ 999.npy}
\end{pipelinebox}


% ---------------------------------------------------------------
\section{소계획 7: GART 7-View 학습}
\label{sec:sp7}

\subsection{현재 상태 vs 목표}
\begin{center}
\renewcommand{\arraystretch}{1.2}
\begin{tabular}{lcc}
\toprule
& \textbf{현재 (6v/100f)} & \textbf{목표 (7v/1000f)} \\
\midrule
카메라 수 & 6 & \textbf{7} \\
프레임 수 & 100 & \textbf{1000} \\
학습 이미지 & 600 & \textbf{7000} \\
학습 스텝 & 3,000 & \textbf{30,000} \\
IMAGE\_ZOOM & 0.5 & \textbf{0.25} (메모리 절약) \\
POSE\_OPT & 비활성 & \textbf{활성 (스텝 15,000부터)} \\
SH 차수 & 1 & \textbf{2} \\
가우시안 초기 수 & 10,475 & \textbf{10,475} \\
\bottomrule
\end{tabular}
\end{center}

\subsection{소계획}
\begin{pipelinebox}[title={SP-7: GART 7-View 30,000스텝 학습}]
\begin{enumerate}
\item GART 코드 수정: \texttt{training\_view=[0,1,2,3,4,5,6]} (1줄)
\item 설정 파일 작성:
  \begin{itemize}
  \item \texttt{TOTAL\_steps: 30000}
  \item \texttt{N\_POSES\_PER\_STEP: 4} (GPU 메모리 최적)
  \item \texttt{IMAGE\_ZOOM\_RATIO: 0.25} (480$\times$270)
  \item \texttt{MAX\_SPH\_ORDER: 2}
  \item \texttt{DENSIFY\_END: 20000}
  \item \texttt{POSE\_OPT\_START: 15000} (전반부: 가우시안 학습, 후반부: 포즈 정제)
  \end{itemize}
\item 학습 실행 (A100, 예상 30--60분)
\item 체크포인트 저장: 매 5000스텝 + 최종
\item PSNR/SSIM 측정: 매 1000스텝 로깅
\end{enumerate}
\textbf{산출물}: \texttt{model.pth}, PSNR/SSIM 곡선, 렌더링 이미지
\end{pipelinebox}


% ---------------------------------------------------------------
\section{소계획 8: GART 포토메트릭 포즈 정제}
\label{sec:sp8}

\subsection{원리}
GART가 학습된 후, 가우시안을 고정하고 SMPL 포즈만 최적화합니다.
포토메트릭 손실은 키포인트(231값)의 \textbf{63,000배} 밀도인 14.5M 픽셀로 감독합니다.

\subsection{소계획}
\begin{pipelinebox}[title={SP-8: 포토메트릭 포즈 정제 루프}]
\begin{enumerate}
\item Phase 1 (Coarse): 가우시안 고정, $\vtheta$ 만 최적화
  \begin{itemize}
  \item lr = $10^{-3}$, 500 iterations
  \item 손실: $\mathcal{L}_{\text{photo}} = (1-\lambda_s)\|I - \hat{I}\|_1 + \lambda_s \mathcal{L}_{\text{SSIM}}$
  \end{itemize}
\item Phase 2 (Fine): 가우시안 + $\vtheta$ 동시 최적화
  \begin{itemize}
  \item 가우시안 lr = 표준, $\vtheta$ lr = $10^{-5}$
  \item 3000 iterations
  \end{itemize}
\item 위상 적응형 시간 정규화:
  \begin{itemize}
  \item 가속기: $\lambda_{\text{temp}} = 0.5$ (7000°/s 허용)
  \item 와인드업: $\lambda_{\text{temp}} = 10.0$ (강한 스무딩)
  \end{itemize}
\item 정제된 SMPL 파라미터 추출 $\to$ 최종 포즈
\end{enumerate}
\textbf{산출물}: 포토메트릭 정제된 1000 프레임 SMPL 파라미터
\end{pipelinebox}


% ---------------------------------------------------------------
\section{소계획 9: SMPL $\to$ 생체역학 브릿지}
\label{sec:sp9}

\subsection{현재 공백}
\begin{itemize}
\item 가상 마커 추출: \textbf{미구현}
\item TRC 파일 생성: \textbf{미구현}
\item OpenSim IK/ID: \textbf{미구현}
\item OpenSim 모델 파일: \textbf{미확보}
\end{itemize}

\subsection{소계획}
\begin{pipelinebox}[title={SP-9: SMPL $\to$ OpenSim 생체역학 브릿지}]
\begin{enumerate}
\item SAM4Dcap의 SMPL $\to$ 105마커 BSM 매핑 코드 적용
  \begin{itemize}
  \item SMPL 6,890 꼭짓점 $\to$ 105 해부학적 마커 (인덱스 매핑)
  \end{itemize}
\item TRC 파일 생성 (OpenSim 표준 형식)
\item 좌표계 변환: 카메라 $\to$ 월드 (R, t, 신장 스케일링)
\item 지면 정렬: $\Delta y = -\min_t(\min_i(p_{y,i}^t))$
\item OpenSim Rajagopal2016 모델로 IK 실행
\item 관절 각도 추출: 어깨 외/내회전, 팔꿈치 굴곡, 체간 회전
\end{enumerate}
\textbf{산출물}: 관절 각도 시계열, TRC 파일, .mot 파일
\end{pipelinebox}


% ---------------------------------------------------------------
\section{소계획 10: Vicon 검증 인프라}
\label{sec:sp10}

\subsection{현재 상태}
\begin{itemize}
\item Vicon C3D: \texttt{set\_002.c3d} (2.9MB), CSV (4.3MB) 존재
\item 42 마커 (Plug-in Gait), 240Hz, 프레임 512--1386 (875프레임)
\item 시간 동기화: \textbf{미구현}
\item 공간 정렬: Procrustes 코드 존재하나 래퍼 없음
\end{itemize}

\subsection{소계획}
\begin{pipelinebox}[title={SP-10: Vicon 검증 파이프라인}]
\begin{enumerate}
\item C3D 파싱: \texttt{c3d} 라이브러리로 42 마커 궤적 추출
\item 시간 동기화: 손목 Y축 교차상관 $\to$ $\Delta t^*$ (±1--2프레임 정밀도)
  \[
  \Delta t^* = \arg\max_{\Delta t} \sum_t s_{\text{RGB}}(t) \cdot s_{\text{Vicon}}(t + \Delta t)
  \]
\item 공간 정렬: $\geq$4 대응점 Procrustes ($s, \mR, \vt$)
\item 프레임 대응 테이블 생성: RGB frame $\leftrightarrow$ Vicon frame
\item MPJPE, PA-MPJPE, 관절각 RMSE, Bland-Altman 계산
\end{enumerate}
\textbf{산출물}: 시간/공간 정렬 매핑, 검증 메트릭 보고서
\end{pipelinebox}


% ---------------------------------------------------------------
\section{소계획 11: 최종 분석 및 시각화}
\label{sec:sp11}

\subsection{소계획}
\begin{pipelinebox}[title={SP-11: 투구 생체역학 결과 분석}]
\begin{enumerate}
\item 투구 이벤트 감지: FC, MER, BR, MIR (기구현 코드 활용)
\item 6단계 위상 분할
\item 주요 지표 추출:
  \begin{itemize}
  \item MER 각도 (°)
  \item 어깨 내회전 속도 (°/s)
  \item 팔꿈치 외반 각도 (°)
  \item 보폭 비율 (\% 신장)
  \item 체간-골반 분리각 (°)
  \end{itemize}
\item GART 렌더링 영상 + 관절각 오버레이 비디오 생성
\item Vicon 대비 Bland-Altman 플롯 생성
\end{enumerate}
\textbf{산출물}: 투구 생체역학 보고서, 비교 플롯, 오버레이 비디오
\end{pipelinebox}


% ====================================================================
% CHAPTER 2: 중계획 (소계획 통합)
% ====================================================================
\chapter{중계획: 소계획 통합}

소계획 11개를 의존성에 따라 4개 중계획으로 통합합니다.

\section{의존성 그래프}

\begin{center}
\begin{tikzpicture}[
  node distance=0.8cm,
  sp/.style={draw, rounded corners, fill=lightblue, minimum width=2.2cm, minimum height=0.6cm, align=center, font=\scriptsize},
  mp/.style={draw, rounded corners, fill=lightorange, minimum width=3cm, minimum height=0.7cm, align=center, font=\small\bfseries},
  arrow/.style={->, thick, >=stealth}
]
% 소계획 노드
\node[sp] (sp1) {SP-1\\캘리브레이션};
\node[sp, right=of sp1] (sp2) {SP-2\\2D 포즈};
\node[sp, right=of sp2] (sp3) {SP-3\\U-HMR 초기값};
\node[sp, below=1.5cm of sp2] (sp4) {SP-4\\7-view reproj};
\node[sp, left=of sp4] (sp5) {SP-5\\이미지 준비};
\node[sp, right=of sp4] (sp6) {SP-6\\SMPL 변환};
\node[sp, below=1.5cm of sp4] (sp7) {SP-7\\GART 학습};
\node[sp, right=of sp7] (sp8) {SP-8\\포즈 정제};
\node[sp, below=1.5cm of sp7] (sp9) {SP-9\\생체역학};
\node[sp, left=of sp9] (sp10) {SP-10\\Vicon 검증};
\node[sp, right=of sp9] (sp11) {SP-11\\최종 분석};

% 중계획 그룹
\node[mp, above=0.3cm of sp2] (mp1) {MP-1: 데이터 준비};
\node[mp, left=1cm of sp7] (mp2) {MP-2: 모션 추정};
\node[mp, below left=0.3cm and -1cm of sp9] (mp3) {MP-3: 3DGS 정제};
\node[mp, below=0.3cm of sp11] (mp4) {MP-4: 검증 \& 분석};

% 화살표
\draw[arrow] (sp1) -- (sp4);
\draw[arrow] (sp2) -- (sp4);
\draw[arrow] (sp3) -- (sp4);
\draw[arrow] (sp1) -- (sp5);
\draw[arrow] (sp4) -- (sp6);
\draw[arrow] (sp5) -- (sp7);
\draw[arrow] (sp6) -- (sp7);
\draw[arrow] (sp7) -- (sp8);
\draw[arrow] (sp8) -- (sp9);
\draw[arrow] (sp9) -- (sp11);
\draw[arrow] (sp10) -- (sp11);
\end{tikzpicture}
\end{center}


% ---------------------------------------------------------------
\section{MP-1: 데이터 준비 (SP-1 + SP-2 + SP-3 + SP-5)}

\begin{keyidea}
\textbf{목표}: 7개 카메라의 캘리브레이션, 2D 키포인트, 이미지를 모두 통합하여 후속 단계의 입력을 완성한다.
\end{keyidea}

\begin{center}
\renewcommand{\arraystretch}{1.2}
\begin{tabular}{llll}
\toprule
\textbf{소계획} & \textbf{입력} & \textbf{출력} & \textbf{병렬 가능} \\
\midrule
SP-1 캘리브레이션 & COLMAP Rd2 & annots\_7view.npy & \checkmark \\
SP-2 2D 포즈 & cam1 회전 이미지 & 7-view keypoints & \checkmark \\
SP-3 U-HMR 초기값 & 기존 결과 & body\_pose init & 재활용 \\
SP-5 이미지 준비 & 원본 + cam1 회전 & ZJU-MoCap 이미지 & \checkmark \\
\bottomrule
\end{tabular}
\end{center}

\begin{importantnote}
SP-1, SP-2, SP-5는 \textbf{병렬 실행} 가능. SP-3은 기존 결과 재활용으로 추가 작업 없음.\\
cam1 이미지 회전이 SP-1, SP-2, SP-5의 공통 전제조건.
\end{importantnote}

\textbf{예상 시간}: 1--2시간 (cam1 OpenPose가 가장 오래 걸림)


% ---------------------------------------------------------------
\section{MP-2: 7-View 모션 추정 (SP-4 + SP-6)}

\begin{keyidea}
\textbf{목표}: 7개 카메라 전체를 활용한 sequential reproj fitting으로 SMPL 포즈 품질을 극대화하고, GART 입력 형식으로 변환한다.
\end{keyidea}

\textbf{핵심 변경사항} (현재 대비):
\begin{enumerate}
\item 카메라: 4 $\to$ \textbf{7} (감독 신호 75\% 증가)
\item 정규화: $\lambda = 0.01 \to 0.001$ (10배 완화)
\item 최적화: 120스텝 $\to$ \textbf{500스텝} (4배 증가)
\item Sequential: Rh, Th만 $\to$ \textbf{Rh, Th, body\_pose 모두} 전달
\item 관절각 제한: 없음 $\to$ \textbf{ROM 클램핑}
\end{enumerate}

\textbf{예상 효과}:
\begin{center}
\renewcommand{\arraystretch}{1.2}
\begin{tabular}{lcc}
\toprule
\textbf{변경} & \textbf{예상 감소} & \textbf{근거} \\
\midrule
$\lambda$ 완화 & $-$50$\sim$100px & 오차 바닥 113px 해제 \\
7-view (vs 4-view) & $-$10$\sim$30px & 깊이 제약 강화 \\
스텝 증가 & $-$10$\sim$20px & 수렴 개선 \\
Sequential body\_pose & $-$5$\sim$15px & 시간 일관성 \\
\midrule
\textbf{총 예상} & \textbf{$-$75$\sim$165px} & $\to$ \textbf{$<$80px 목표} \\
\bottomrule
\end{tabular}
\end{center}

\textbf{예상 시간}: 10--20분 (A100, 1000 프레임)


% ---------------------------------------------------------------
\section{MP-3: 3DGS 학습 + 포즈 정제 (SP-7 + SP-8)}

\begin{keyidea}
\textbf{목표}: GART 7-view 30,000스텝 학습 후, 포토메트릭 손실로 SMPL 포즈를 추가 정제하여 키포인트 기반 방법의 한계를 극복한다.
\end{keyidea}

\textbf{2단계 실행}:
\begin{enumerate}
\item \textbf{SP-7}: GART 학습 (30,000스텝, 30--60분)
  \begin{itemize}
  \item 입력: MP-2 출력 (정제된 SMPL) + MP-1 출력 (7-view 이미지/캘리브레이션)
  \item 전반부 (0--15,000): 가우시안 학습만 (SMPL 포즈 고정)
  \item 후반부 (15,000--30,000): 포즈 공동 최적화 시작
  \end{itemize}
\item \textbf{SP-8}: 포토메트릭 정제 (3,500 iterations, 10--20분)
  \begin{itemize}
  \item 가우시안 고정, $\vtheta$ 만 최적화
  \item 위상 적응형 시간 정규화
  \item 14.5M 픽셀/프레임 감독 (키포인트의 63,000배)
  \end{itemize}
\end{enumerate}

\begin{importantnote}
SP-8이 이 프로젝트의 \textbf{핵심 기여점}입니다. 기존 SOTA (Theia3D 52mm, KinaTrax 57mm)는 모두 키포인트 기반이며, 포토메트릭 정제를 적용한 투구 분석은 세계 최초입니다.
\end{importantnote}

\textbf{예상 시간}: 40--80분


% ---------------------------------------------------------------
\section{MP-4: 생체역학 검증 \& 분석 (SP-9 + SP-10 + SP-11)}

\begin{keyidea}
\textbf{목표}: 정제된 SMPL $\to$ OpenSim 관절 각도 추출 $\to$ Vicon 검증 $\to$ 임상 지표 산출.
\end{keyidea}

\textbf{3단계 순차 실행}:
\begin{enumerate}
\item \textbf{SP-9}: SMPL $\to$ TRC $\to$ OpenSim IK (1--2시간)
\item \textbf{SP-10}: Vicon 시간/공간 정렬 (30분)
\item \textbf{SP-11}: 메트릭 계산 + 시각화 (1시간)
\end{enumerate}

\begin{warningbox}
SP-10 (Vicon 동기화)이 실패하면 SP-11의 정량적 검증이 불가능합니다. 이 경우 정성적 평가(렌더링 품질, 관절각 프로파일 형태)만 수행합니다.
\end{warningbox}


% ====================================================================
% CHAPTER 3: 대계획
% ====================================================================
\chapter{대계획: 중계획 통합}

\section{실행 순서 및 의존성}

\begin{center}
\begin{tikzpicture}[
  node distance=2cm,
  mp/.style={draw, rounded corners, fill=lightorange, minimum width=4cm, minimum height=1cm, align=center, font=\bfseries},
  arrow/.style={->, very thick, >=stealth}
]
\node[mp] (mp1) {MP-1: 데이터 준비\\(1--2시간)};
\node[mp, below=of mp1] (mp2) {MP-2: 7-View 모션 추정\\(10--20분)};
\node[mp, below=of mp2] (mp3) {MP-3: 3DGS 학습 + 정제\\(40--80분)};
\node[mp, below=of mp3] (mp4) {MP-4: 검증 \& 분석\\(2--4시간)};

\draw[arrow] (mp1) -- (mp2) node[midway, right, font=\small] {7-view K,R,t + keypoints + images};
\draw[arrow] (mp2) -- (mp3) node[midway, right, font=\small] {정제된 SMPL params ($<$80px)};
\draw[arrow] (mp3) -- (mp4) node[midway, right, font=\small] {포토메트릭 정제 SMPL ($<$20px 목표)};
\end{tikzpicture}
\end{center}

\section{리스크 분석}

\begin{center}
\renewcommand{\arraystretch}{1.3}
\begin{tabular}{p{4cm}ccp{5cm}}
\toprule
\textbf{리스크} & \textbf{확률} & \textbf{영향} & \textbf{대응} \\
\midrule
cam1 회전 후 OpenPose 실패 & 낮음 & 낮음 & 6-view로 진행 (cam1 제외) \\
7-view reproj이 4-view보다 나빠짐 & 낮음 & 중간 & $\lambda$ 재조정, 4-view 폴백 \\
GART A100 OOM & 중간 & 중간 & ZOOM=0.125, N\_POSES=2 \\
GART 포즈 정제 발산 & 중간 & 높음 & lr 감소, 위상 마스킹 \\
Vicon 동기화 실패 & 중간 & 중간 & 정성적 평가만 수행 \\
OpenSim 모델 미확보 & 낮음 & 높음 & Rajagopal2016 범용 모델 \\
\bottomrule
\end{tabular}
\end{center}

\section{품질 게이트}

각 중계획 완료 시 다음 단계 진행 여부를 판단합니다:

\begin{center}
\renewcommand{\arraystretch}{1.3}
\begin{tabular}{lll}
\toprule
\textbf{게이트} & \textbf{통과 기준} & \textbf{실패 시} \\
\midrule
G1 (MP-1 $\to$ MP-2) & 7개 카메라 K,R,t + keypoints 완비 & cam1 제외 6-view로 진행 \\
G2 (MP-2 $\to$ MP-3) & reproj $<$ 100px (7-view 평균) & $\lambda$ 재조정 후 재실행 \\
G3 (MP-3 $\to$ MP-4) & PSNR $>$ 20dB, 시각적 형상 확인 & 스텝 증가 (30K $\to$ 60K) \\
\bottomrule
\end{tabular}
\end{center}


% ====================================================================
% CHAPTER 4: 전체 계획
% ====================================================================
\chapter{전체 계획: 최종 통합}

\section{한 눈에 보는 전체 흐름}

\begin{center}
\begin{tikzpicture}[
  node distance=0.6cm,
  planstep/.style={draw, rounded corners, fill=lightblue, minimum width=12cm, minimum height=0.7cm, align=left, font=\small},
  gate/.style={draw, diamond, fill=lightorange, minimum width=1.5cm, minimum height=0.5cm, font=\scriptsize, align=center},
  arrow/.style={->, thick, >=stealth}
]

\node[planstep] (s1) {1. cam1 이미지 90° 회전 (1000장)};
\node[planstep, below=of s1] (s2) {2. cam1 OpenPose 실행 + 4단계 스무딩};
\node[planstep, below=of s2] (s3) {3. COLMAP Rd2 $\to$ 7-view annots.npy 변환 (cam1 K 보정 포함)};
\node[planstep, below=of s3] (s4) {4. ZJU-MoCap 이미지 디렉토리 구성 (7$\times$1000 = 7000장)};

\node[gate, below=0.8cm of s4] (g1) {G1};

\node[planstep, below=0.8cm of g1] (s5) {5. 7-view sequential reproj fitting ($\lambda$=0.001, 500스텝, ROM 클램핑)};
\node[planstep, below=of s5] (s6) {6. SMPL params $\to$ smpl\_params/*.npy 변환};

\node[gate, below=0.8cm of s6] (g2) {G2\\$<$100px?};

\node[planstep, below=0.8cm of g2] (s7) {7. GART 7-view 학습 (30,000스텝, ZOOM=0.25, SH=2)};
\node[planstep, below=of s7] (s8) {8. 포토메트릭 포즈 정제 (가우시안 고정, $\vtheta$ lr=$10^{-3}$)};

\node[gate, below=0.8cm of s8] (g3) {G3\\PSNR$>$20?};

\node[planstep, below=0.8cm of g3] (s9) {9. SMPL $\to$ 105마커 $\to$ TRC $\to$ OpenSim IK};
\node[planstep, below=of s9] (s10) {10. Vicon C3D 파싱 + 시간/공간 동기화};
\node[planstep, below=of s10] (s11) {11. MPJPE, 관절각 RMSE, Bland-Altman, 투구 이벤트 분석};

\draw[arrow] (s1) -- (s2);
\draw[arrow] (s2) -- (s3);
\draw[arrow] (s3) -- (s4);
\draw[arrow] (s4) -- (g1);
\draw[arrow] (g1) -- (s5);
\draw[arrow] (s5) -- (s6);
\draw[arrow] (s6) -- (g2);
\draw[arrow] (g2) -- (s7);
\draw[arrow] (s7) -- (s8);
\draw[arrow] (s8) -- (g3);
\draw[arrow] (g3) -- (s9);
\draw[arrow] (s9) -- (s10);
\draw[arrow] (s10) -- (s11);
\end{tikzpicture}
\end{center}


\section{총 예상 시간}

\begin{center}
\renewcommand{\arraystretch}{1.3}
\begin{tabular}{llc}
\toprule
\textbf{단계} & \textbf{내용} & \textbf{시간} \\
\midrule
MP-1 & cam1 회전 + OpenPose + 캘리브레이션 변환 + 이미지 구조화 & 1--2시간 \\
MP-2 & 7-view reproj fitting + SMPL 변환 & 10--20분 \\
MP-3 & GART 30K스텝 + 포토메트릭 정제 & 40--80분 \\
MP-4 & OpenSim + Vicon 검증 + 분석 & 2--4시간 \\
\midrule
\textbf{총계} & & \textbf{4--8시간} \\
\bottomrule
\end{tabular}
\end{center}


\section{예상 최종 성과}

\begin{center}
\renewcommand{\arraystretch}{1.3}
\begin{tabular}{lccc}
\toprule
\textbf{지표} & \textbf{현재} & \textbf{예상 (reproj 후)} & \textbf{예상 (GART 정제 후)} \\
\midrule
Reproj 오차 (px) & 153.5 & $<$80 & $<$20 \\
MPJPE (mm) & 미측정 & $<$60 & $<$30 \\
어깨 외회전 RMSE (°) & 미측정 & $<$15 & $<$8 \\
팔꿈치 굴곡 RMSE (°) & 미측정 & $<$10 & $<$5 \\
렌더링 PSNR (dB) & -- & -- & $>$22 \\
\bottomrule
\end{tabular}
\end{center}


\section{산출물 체크리스트}

\begin{checklistbox}[title={최종 산출물}]
\begin{enumerate}
\item[$\square$] 7-view annots.npy (캘리브레이션)
\item[$\square$] 7-view 스무딩 OpenPose 키포인트 (7$\times$1000 프레임)
\item[$\square$] ZJU-MoCap 이미지 디렉토리 (7000 이미지)
\item[$\square$] 7-view reproj 정제 SMPL 파라미터 (1000 프레임)
\item[$\square$] GART 학습 모델 (model.pth, 30K스텝)
\item[$\square$] 포토메트릭 정제 SMPL 파라미터 (1000 프레임)
\item[$\square$] PSNR/SSIM 수렴 곡선
\item[$\square$] 7-view 렌더링 vs GT 비교 이미지/비디오
\item[$\square$] TRC 마커 궤적 파일
\item[$\square$] OpenSim IK 관절 각도 (.mot)
\item[$\square$] Vicon 동기화 매핑
\item[$\square$] MPJPE, 관절각 RMSE 보고서
\item[$\square$] Bland-Altman 플롯
\item[$\square$] 투구 이벤트 (FC/MER/BR/MIR) + 위상 분할
\item[$\square$] 임상 지표 (MER°, 내회전속도, 보폭비율)
\end{enumerate}
\end{checklistbox}


\begin{summary}
\textbf{전체 계획 요약}:
\begin{enumerate}
\item \textbf{MP-1 (데이터)}: cam1 세로$\to$가로 회전, 7-view 캘리브레이션/키포인트/이미지 통합
\item \textbf{MP-2 (모션)}: 정규화 완화($\lambda$=0.001) + 7-view + 500스텝 + ROM 클램핑 $\to$ reproj $<$80px
\item \textbf{MP-3 (3DGS)}: GART 30K스텝 학습 + 포토메트릭 포즈 정제 $\to$ 14.5M 픽셀 감독
\item \textbf{MP-4 (검증)}: SMPL$\to$TRC$\to$OpenSim IK + Vicon 검증 $\to$ MPJPE, 관절각 RMSE
\end{enumerate}
11개 소계획, 4개 중계획, 3개 품질 게이트, 총 4--8시간.\\
\textbf{핵심 기여}: 포토메트릭 3DGS 포즈 정제를 투구 분석에 적용한 세계 최초 시스템.
\end{summary}


% ====================================================================
% CHAPTER 5: 세부 구현 명세 (Claude 스킬 형식)
% ====================================================================
\chapter{세부 구현 명세 --- Claude 실행 스킬}

이 장은 각 소계획의 세부 구현을 \textbf{Claude가 서버에서 바로 실행할 수 있는 스킬 형식}으로 작성합니다. 각 스킬은 서버 경로, 정확한 코드, 실행 명령, 검증 기준을 포함합니다.


% ---------------------------------------------------------------
\section{SKILL-1: cam1 이미지 회전}

\begin{codebox}[title={서버 경로 및 환경}]
\begin{verbatim}
서버: /home/elicer/
conda env: base (또는 아무 환경)
입력: /home/elicer/images/cam1/ (1080x1920 세로 JPEG, 1000장)
출력: /home/elicer/images_rotated/cam1/ (1920x1080 가로 JPEG)
\end{verbatim}
\end{codebox}

\begin{codebox}[title={Python 코드}]
\begin{verbatim}
import cv2, os, glob
src = '/home/elicer/images/cam1'
dst = '/home/elicer/images_rotated/cam1'
os.makedirs(dst, exist_ok=True)
for f in sorted(glob.glob(f'{src}/*.jpg')):
    img = cv2.imread(f)           # (1920, 1080, 3)
    rot = cv2.rotate(img, cv2.ROTATE_90_CLOCKWISE)  # (1080, 1920, 3)
    cv2.imwrite(f'{dst}/{os.path.basename(f)}', rot)
print(f'Rotated {len(os.listdir(dst))} images')
\end{verbatim}
\end{codebox}

\begin{codebox}[title={cam1 K행렬 변환 (COLMAP SIMPLE\_RADIAL)}]
\begin{verbatim}
# COLMAP cam1 원본 (세로): f=733.18, cx=540.0, cy=960.0, k1=-0.00282
# 90도 CW 회전 후 (가로):
#   K_rot = [[733.18, 0, 960.0],   ← cx_new = H_orig - cy_orig = 1920 - 960
#            [0, 733.18, 540.0],   ← cy_new = cx_orig = 540
#            [0, 0, 1]]
# R_rot = R_orig @ [[0, -1, 0], [1, 0, 0], [0, 0, 1]]
\end{verbatim}
\end{codebox}

\textbf{검증}: 회전된 이미지가 1920$\times$1080이고 투구자가 정상 방향인지 확인.


% ---------------------------------------------------------------
\section{SKILL-2: cam1 OpenPose 실행}

\begin{codebox}[title={서버 명령어}]
\begin{verbatim}
# cam1 회전 이미지에서 OpenPose 실행
cd /home/elicer/openpose
./build/examples/openpose/openpose.bin \
    --image_dir /home/elicer/images_rotated/cam1 \
    --write_json /home/elicer/openpose_cam1_rot/ \
    --display 0 --render_pose 0 \
    --model_pose BODY_25
\end{verbatim}
\end{codebox}

\begin{codebox}[title={스무딩 (기존과 동일 파라미터)}]
\begin{verbatim}
from scipy.signal import savgol_filter
window = 11; poly = 3  # openpose_smoothed_v2와 동일
# 1000 프레임 x 25 관절 x 3 (x, y, conf)에 적용
# x, y만 스무딩, conf는 유지
\end{verbatim}
\end{codebox}

\textbf{검증}: \texttt{openpose\_smoothed\_v2/cam1/} 에 1000개 JSON 파일 생성 확인.


% ---------------------------------------------------------------
\section{SKILL-3: 7-View 캘리브레이션 통합}

\begin{codebox}[title={COLMAP $\to$ annots.npy 변환}]
\begin{verbatim}
import numpy as np, json

# COLMAP Round 2 파라미터 (SIMPLE_RADIAL)
cams = {
  'cam1': {'f': 733.18, 'cx': 960.0, 'cy': 540.0, 'k1': -0.00282,
           'w': 1920, 'h': 1080},  # 회전 후!
  'cam2': {'f': 744.46, 'cx': 960.0, 'cy': 540.0, 'k1': -0.00602,
           'w': 1920, 'h': 1080},
  'cam3': {'f': 740.98, 'cx': 960.0, 'cy': 540.0, 'k1':  0.01492,
           'w': 1920, 'h': 1080},
  'cam4': {'f': 530.77, 'cx': 960.0, 'cy': 540.0, 'k1': -0.00362,
           'w': 1920, 'h': 1080},
  'cam5': {'f': 586.39, 'cx': 960.0, 'cy': 540.0, 'k1': -0.06382,
           'w': 1920, 'h': 1080},
  'cam6': {'f':1367.18, 'cx': 960.0, 'cy': 540.0, 'k1': -0.34613,
           'w': 1920, 'h': 1080},  # 주의: k1 매우 큼!
  'cam7': {'f': 811.28, 'cx': 960.0, 'cy': 540.0, 'k1': -0.10493,
           'w': 1920, 'h': 1080},
}

# K 행렬 생성 (SIMPLE_RADIAL: fx=fy=f)
for name, c in cams.items():
    c['K'] = np.array([[c['f'], 0, c['cx']],
                       [0, c['f'], c['cy']],
                       [0, 0, 1]])
    c['D'] = np.array([c['k1'], 0, 0, 0, 0])  # k1만 사용

# R, t는 COLMAP binary에서 추출 (read_model.py)
# 또는 VGGT JSON에서 로드
\end{verbatim}
\end{codebox}

\begin{warningbox}
\textbf{cam6 왜곡 주의}: $k_1 = -0.346$은 매우 큰 핀쿠션 왜곡입니다.
왜곡 보정 없이 reproj하면 cam6에서 50--100px 추가 오차가 발생합니다.\\
\textbf{해결}: reproj 전에 \texttt{cv2.undistortPoints()}로 2D 키포인트를 보정하거나,\\
투영 시 왜곡 모델을 포함합니다.
\end{warningbox}

\begin{importantnote}
\textbf{좌표계 문제}: 현재 코드는 VGGT 캘리브레이션 + \texttt{R\_fix\_vggt\_to\_smpl.npy}를 사용합니다.
COLMAP으로 전환 시 새로운 좌표 정렬이 필요합니다.\\
\textbf{권장}: 기존 VGGT 캘리브레이션을 유지하되, cam1만 추가합니다.
VGGT JSON에 cam1이 이미 포함되어 있으므로 (vggt\_calibrate.py에서 처리),
추가 캘리브레이션 없이 cam1을 로드할 수 있습니다.
\end{importantnote}


% ---------------------------------------------------------------
\section{SKILL-4: 7-View Sequential Reproj Fitting (핵심)}

이것이 전체 파이프라인에서 가장 중요한 스크립트입니다.

\begin{codebox}[title={주요 변경사항 vs 현재 uhmr\_7cam\_reproj.py}]
\begin{verbatim}
### 변경 1: 카메라 선택 (line 18)
# 현재: 7개 로드하지만 실제로 4개만 사용
# 변경: 7개 전체 사용 (또는 cam6 왜곡 문제시 6개)

### 변경 2: 배치 → 진짜 Sequential (구조 변경)
# 현재: BATCH=50, 50프레임씩 동일 초기값
# 변경: frame-by-frame, prev_Rh/Th/bp → 다음 프레임 init

### 변경 3: 정규화 (line 146)
# 현재: loss_reg = 0.01 * ((bp_tune - bp_fixed)**2).sum() / B
# 변경: loss_reg = 0.001 * ((bp_tune - bp_fixed)**2).sum()

### 변경 4: Stage 2 강화 (lines 137-151)
# 현재: 100스텝, lr=0.005
# 변경: 300스텝, lr=0.01

### 변경 5: 관절각 클램핑 (새로 추가)
# 매 스텝 후:
#   elbow_idx = [18*3:18*3+3, 19*3:19*3+3]  # SMPL left/right elbow
#   bp_tune[elbow_range].clamp_(-2.6, 0)  # 0~150도 굴곡

### 변경 6: 왜곡 보정 (cam6 필수)
# OpenPose 2D keypoint를 undistort:
#   pts_undist = cv2.undistortPoints(pts, K, D, P=K)
# 또는 투영 시 radial distortion 적용

### 변경 7: GMoF sigma (line 81)
# 현재: sigma=100.0
# 변경: sigma=50.0 (이상치 더 강하게 억제)
\end{verbatim}
\end{codebox}

\begin{codebox}[title={핵심 최적화 루프 (pseudo-code)}]
\begin{verbatim}
prev_Rh = go_init.copy()
prev_Th = Th_init.copy()
prev_bp = body_pose_init.clone()

for fi in range(1000):
    # 초기화: 이전 프레임 결과
    Rh = prev_Rh.clone().requires_grad_(True)
    Th = prev_Th.clone().requires_grad_(True)
    bp = prev_bp.clone().requires_grad_(True)

    kp2d = kp2d_all[fi]   # (7, 25, 2)
    conf = conf_all[fi]   # (7, 25)

    # Stage 1: RT only (200 steps)
    opt1 = Adam([Rh, Th], lr=0.03)
    for s in range(200):
        out = smpl(global_orient=Rh, body_pose=bp, betas=betas, transl=Th)
        j25 = J_reg @ out.vertices
        projs = project_7cam(j25, K_all, R_all, t_all)  # (7, 25, 2)
        loss = (conf[...,None] * gmof(projs - kp2d, sigma=50)).sum()
        opt1.zero_grad(); loss.backward(); opt1.step()

    # Stage 2: + body_pose (300 steps)
    opt2 = Adam([Rh, Th, bp], lr=0.01)
    for s in range(300):
        out = smpl(global_orient=Rh, body_pose=bp, betas=betas, transl=Th)
        j25 = J_reg @ out.vertices
        projs = project_7cam(j25, K_all, R_all, t_all)
        loss_r = (conf[...,None] * gmof(projs - kp2d, sigma=50)).sum()
        loss_reg = 0.001 * ((bp - body_pose_init)**2).sum()
        loss = loss_r + loss_reg
        opt2.zero_grad(); loss.backward(); opt2.step()

        # 관절각 클램핑
        with torch.no_grad():
            bp[18*3:18*3+3].clamp_(-2.6, 0)  # R elbow
            bp[19*3:19*3+3].clamp_(-2.6, 0)  # L elbow

    prev_Rh = Rh.detach().clone()
    prev_Th = Th.detach().clone()
    prev_bp = bp.detach().clone()

    # 저장 + 로깅
    save_frame(fi, Rh, Th, bp, betas)
    reproj_err = compute_reproj_error(fi)
    print(f'Frame {fi}: reproj={reproj_err:.1f}px')
\end{verbatim}
\end{codebox}

\begin{warningbox}
\textbf{frame-by-frame는 batch보다 느림}: 1000프레임 $\times$ 500스텝 = 50만 iterations.\\
A100에서 약 20--40분 예상 (batch 방식의 4--5분 대비 느리지만 품질이 핵심).\\
\textbf{대안}: mini-batch=10으로 sequential 효과 유지하면서 속도 개선 가능.
\end{warningbox}

\textbf{검증 기준}:
\begin{itemize}
\item 평균 reproj $<$ 80px (현재 153.5px에서 48\% 이상 감소)
\item 95\% 프레임 $<$ 100px
\item 최악 프레임 $<$ 150px
\end{itemize}


% ---------------------------------------------------------------
\section{SKILL-5: ZJU-MoCap 데이터 구조화}

\begin{codebox}[title={디렉토리 구조 생성}]
\begin{verbatim}
base = '/home/elicer/gart_7view_data'
# 이미지
mkdir -p $base/images/{0,1,2,3,4,5,6}
# cam1(회전) → 0/, cam2 → 1/, ..., cam7 → 6/
ln -s /home/elicer/images_rotated/cam1/*.jpg $base/images/0/
ln -s /home/elicer/images/cam2/*.jpg $base/images/1/
ln -s /home/elicer/images/cam3/*.jpg $base/images/2/
ln -s /home/elicer/images/cam4/*.jpg $base/images/3/
ln -s /home/elicer/images/cam5/*.jpg $base/images/4/
ln -s /home/elicer/images/cam6/*.jpg $base/images/5/
ln -s /home/elicer/images/cam7/*.jpg $base/images/6/
# SMPL params
mkdir -p $base/smpl_params
\end{verbatim}
\end{codebox}


% ---------------------------------------------------------------
\section{SKILL-6: SMPL 파라미터 변환}

\begin{codebox}[title={EasyMocap JSON $\to$ GART NPY}]
\begin{verbatim}
import numpy as np, json, os
src = '/home/elicer/reproj_7view/smpl'      # SKILL-4 출력
dst = '/home/elicer/gart_7view_data/smpl_params'
os.makedirs(dst, exist_ok=True)

for fi in range(1000):
    with open(f'{src}/{fi:06d}.json') as f:
        d = json.load(f)[0]
    np.save(f'{dst}/{fi:03d}.npy', {
        'global_orient': np.array(d['Rh'][0]),      # (3,)
        'body_pose': np.array(d['poses'][0])[3:],   # (69,) skip global
        'transl': np.array(d['Th'][0]),              # (3,)
        'betas': np.array(d['shapes'][0]),           # (10,)
    })
\end{verbatim}
\end{codebox}


% ---------------------------------------------------------------
\section{SKILL-7: GART 7-View 학습}

\begin{codebox}[title={GART 코드 수정 (1줄)}]
\begin{verbatim}
# 파일: $GART_ROOT/lib_data/zju_mocap.py, line 81
# 변경 전: training_view=[4],
# 변경 후:
training_view=[0, 1, 2, 3, 4, 5, 6],
\end{verbatim}
\end{codebox}

\begin{codebox}[title={학습 실행}]
\begin{verbatim}
conda activate gart
cd /home/elicer/GART

python train.py \
  --data_dir /home/elicer/gart_7view_data \
  --total_steps 30000 \
  --n_poses_per_step 4 \
  --image_zoom_ratio 0.25 \
  --max_sph_order 2 \
  --densify_end 20000 \
  --pose_opt_start 15000 \
  --save_interval 5000 \
  --eval_interval 1000
\end{verbatim}
\end{codebox}

\textbf{검증 기준}: PSNR $>$ 20dB (스텝 30,000), 시각적으로 투구자 형상 선명.


% ---------------------------------------------------------------
\section{SKILL-8: 포토메트릭 포즈 정제}

\begin{codebox}[title={GART 내장 포즈 최적화 활용}]
\begin{verbatim}
# POSE_OPT_START=15000 이면 스텝 15K부터 자동 활성화
# 추가 정제가 필요한 경우:

python refine_pose.py \
  --model_path /home/elicer/gart_7view_data/model_30000.pth \
  --data_dir /home/elicer/gart_7view_data \
  --pose_lr 1e-3 \
  --fix_gaussians \
  --iterations 500 \
  --phase_adaptive \
  --lambda_temp_accel 0.5 \
  --lambda_temp_windup 10.0
\end{verbatim}
\end{codebox}

\begin{importantnote}
\texttt{refine\_pose.py}는 \textbf{별도 작성 필요}. GART의 기존 학습 루프를 수정하여:
\begin{enumerate}
\item 가우시안 파라미터 \texttt{requires\_grad=False}
\item SMPL $\vtheta$만 \texttt{requires\_grad=True}
\item FC/MER/BR/MIR 이벤트로 위상 감지 $\to$ $\lambda_{\text{temp}}$ 적응
\end{enumerate}
\end{importantnote}


% ---------------------------------------------------------------
\section{SKILL-9: SMPL $\to$ TRC 변환}

\begin{codebox}[title={105마커 BSM 매핑 (SAM4Dcap 참조)}]
\begin{verbatim}
import numpy as np, smplx

model = smplx.create(model_path, model_type='smpl', gender='neutral')

# 105 BSM 마커 → SMPL 꼭짓점 인덱스 매핑
# (SAM4Dcap의 smpl2addbm/marker_indices_105.json 참조)
# 핵심 마커 예시:
MARKER_MAP = {
    'RSHO': 3011,  # Right shoulder
    'RELB': 1620,  # Right elbow
    'RWRA': 2098,  # Right wrist A
    'RKNE': 1002,  # Right knee
    # ... 105개 전체
}

def smpl_to_trc(smpl_params, output_path, fps=240):
    """SMPL params → TRC file for OpenSim."""
    with open(output_path, 'w') as f:
        # TRC header
        f.write('PathFileType\t4\t(X/Y/Z)\tmarkers.trc\n')
        f.write(f'DataRate\tCameraRate\tNumFrames\tNumMarkers\t...\n')
        f.write(f'{fps}\t{fps}\t{len(smpl_params)}\t{len(MARKER_MAP)}\t...\n')
        # Column headers
        # Data rows: frame, time, marker_x, marker_y, marker_z, ...
        for fi, params in enumerate(smpl_params):
            out = model(body_pose=params['body_pose'], ...)
            vertices = out.vertices[0].numpy()
            markers = {name: vertices[idx] for name, idx in MARKER_MAP.items()}
            # Write row
\end{verbatim}
\end{codebox}


% ---------------------------------------------------------------
\section{SKILL-10: Vicon 시간/공간 동기화}

\begin{codebox}[title={C3D 파싱 + 교차상관 동기화}]
\begin{verbatim}
import c3d, numpy as np
from scipy.signal import correlate

# 1. C3D 파싱
reader = c3d.Reader(open('/home/elicer/data/markerbased/subject_1/'
                          'set_002/set_002.c3d', 'rb'))
vicon_markers = {}
for i, pts, analog in reader.read_frames():
    # pts: (N_markers, 5) = x, y, z, residual, camera_count
    vicon_markers[i] = pts[:, :3]  # mm

# 2. 교차상관 시간 동기화
# RGB 신호: 오른쪽 손목 Y축 (OpenPose joint 4)
signal_rgb = openpose_3d[:, 4, 1]  # (1000,)
# Vicon 신호: RWRA marker Y축
signal_vicon = vicon_rwra[:, 1]    # (875,)

# 정규화
s1 = (signal_rgb - signal_rgb.mean()) / signal_rgb.std()
s2 = (signal_vicon - signal_vicon.mean()) / signal_vicon.std()

# 교차상관
corr = correlate(s1, s2, mode='full')
lags = np.arange(-len(s2)+1, len(s1))
offset = lags[np.argmax(np.abs(corr))]
print(f'Time offset: {offset} frames ({offset/240:.3f}s)')

# 3. 공간 정렬 (Procrustes)
# 4+ 대응점으로 s, R, t 추정
\end{verbatim}
\end{codebox}

\textbf{Vicon 데이터 위치}:
\begin{itemize}
\item C3D: \texttt{/data/markerbased/subject\_1/set\_002/set\_002.c3d} (2.9MB)
\item CSV: \texttt{set\_002.csv} (4.3MB) --- 궤적은 line 5278--6158
\item 프레임 범위: 512--1386 (875 프레임, 240Hz)
\item 마커 수: 42 (Plug-in Gait)
\item 힘판: 2$\times$AMTI @ 1200Hz
\end{itemize}


% ---------------------------------------------------------------
\section{SKILL-11: 최종 분석}

\begin{codebox}[title={투구 이벤트 감지 + 메트릭}]
\begin{verbatim}
from pitching_pipeline.src.biomechanics.pitching_events import (
    detect_pitching_events
)
from pitching_pipeline.src.validation.metrics import compute_all_metrics

# 투구 이벤트 감지
events = detect_pitching_events(
    joints_3d,          # (1000, 25, 3)
    fps=240.0,
    throwing_side="right"
)
# events = {'foot_contact': 142, 'max_external_rotation': 398,
#           'ball_release': 405, 'max_internal_rotation': 430}

# 검증 메트릭
metrics = compute_all_metrics(
    pred_joints=smpl_joints,   # (N, 24, 3)
    gt_joints=vicon_joints,    # (N, 24, 3) (동기화 후)
)
# metrics = {'mpjpe_mm': 28.3, 'pa_mpjpe_mm': 22.1, ...}
\end{verbatim}
\end{codebox}

\begin{codebox}[title={검증 대상 관절 인덱스 (OpenPose Body25)}]
\begin{verbatim}
# 투구 분석 핵심 관절:
RIGHT_SHOULDER = 2   # RSho
RIGHT_ELBOW    = 3   # RElb
RIGHT_WRIST    = 4   # RWri
NECK           = 1   # Neck (체간 기준)
MID_HIP        = 8   # MidHip (골반 기준)
RIGHT_ANKLE    = 11  # RAnk (보폭 계산)
LEFT_ANKLE     = 14  # LAnk (보폭 계산)
\end{verbatim}
\end{codebox}

\textbf{검증 목표}:
\begin{center}
\renewcommand{\arraystretch}{1.2}
\begin{tabular}{lcc}
\toprule
\textbf{지표} & \textbf{목표 (Excellent)} & \textbf{목표 (Acceptable)} \\
\midrule
MPJPE & $<$20mm & $<$40mm \\
어깨 RMSE & $<$5° & $<$10° \\
팔꿈치 RMSE & $<$3° & $<$7° \\
토크 nRMSE & $<$10\% & $<$20\% \\
\bottomrule
\end{tabular}
\end{center}


% ====================================================================
% CHAPTER 6: 문제점 로그 + 알려진 함정
% ====================================================================
\chapter{알려진 함정 및 문제점 로그}

실험 로그에서 발견된 모든 함정을 정리합니다. Claude가 서버 작업 시 이 목록을 반드시 참조해야 합니다.

\section{추가 발견 (로그 정밀 분석)}

\begin{failbox}
\textbf{F-0 (Critical): VGGT vs COLMAP 캘리브레이션 선택}\\
\textbf{기존 성공한 4-view 파이프라인(153.5px)은 VGGT 캘리브레이션을 사용했습니다.}\\
COLMAP Rd2의 초점거리는 531--1367px (2.6배 편차)로 카메라 간 불일치가 심합니다.\\
VGGT는 $\sim$734px으로 일관적이며 다시점 동시 추정 장점이 있습니다.\\
\textbf{결론}: 7-view 확장 시에도 \textbf{VGGT를 기본}으로 사용하고, R\_fix 변환을 유지합니다.\\
COLMAP은 왜곡 파라미터(k1)만 참고합니다.
\end{failbox}

\begin{failbox}
\textbf{F-0b: $\lambda_\theta$ 권장값 0.001이 아닌 0.0001}\\
fitting\_guide.tex의 분석에 따르면, $\lambda_\theta = 0.001$에서도 오차 바닥 113px가 유지되었습니다.\\
투구 동작(7000°/s)에서는 \textbf{$\lambda_\theta = 0.0001$}이 필요합니다.\\
02\_reproj 스크립트를 $\lambda = 0.0001$로 수정하였습니다.
\end{failbox}

\begin{failbox}
\textbf{F-0c: SMPL 모델 파일 버전}\\
서버에 두 개의 SMPL\_NEUTRAL이 있습니다:\\
--- \texttt{~/models/smpl/SMPL\_NEUTRAL.pkl} (304KB, \textbf{경량 --- 부정확!})\\
--- \texttt{~/EasyMocap/data/smplx/.../SMPL\_NEUTRAL\_FIXED.pkl} (21MB, \textbf{전체})\\
\textbf{반드시 21MB 버전을 사용}합니다 (smplx.create의 model\_path로 EasyMocap 경로 지정).
\end{failbox}

\begin{failbox}
\textbf{F-0d: Vicon 동기화 결과 (로컬에서 완료)}\\
교차상관 결과: \textbf{오프셋 125프레임 (0.52초)}, 상관계수 0.78, 겹침 750프레임.\\
RGB 프레임 0--749 $\leftrightarrow$ Vicon 프레임 125--874.\\
\texttt{data/vicon\_sync/sync\_result.json}에 저장됨.
\end{failbox}

\section{캘리브레이션 함정}

\begin{failbox}
\textbf{F-1: VGGT vs COLMAP 좌표계 불일치}\\
VGGT는 자체 월드 좌표를 생성하고, COLMAP은 다른 월드 좌표를 사용합니다.
현재 코드는 \texttt{R\_fix\_vggt\_to\_smpl.npy}로 VGGT$\to$SMPL 변환을 적용합니다.
COLMAP 캘리브레이션으로 전환하면 이 R\_fix가 맞지 않을 수 있습니다.\\
\textbf{해결}: VGGT 캘리브레이션을 유지하거나, COLMAP 전환 시 새 R\_fix를 계산.
\end{failbox}

\begin{failbox}
\textbf{F-2: cam6 왜곡 $k_1 = -0.346$}\\
cam6는 모든 카메라 중 왜곡이 가장 큽니다. 왜곡 보정 없이 reproj하면 외곽 키포인트에서 50--100px 오차가 추가됩니다.\\
\textbf{해결}: \texttt{cv2.undistortPoints()} 적용 또는 투영 함수에 방사왜곡 모델 추가.
\end{failbox}

\begin{failbox}
\textbf{F-3: cam1 K행렬 회전 변환}\\
세로 이미지를 CW 90° 회전할 때 K행렬의 cx, cy가 바뀝니다.\\
$K_{\text{rot}} = \begin{pmatrix} f & 0 & H_{\text{orig}} - c_y \\ 0 & f & c_x \\ 0 & 0 & 1 \end{pmatrix}$ (단순 교환이 아님!)\\
\textbf{해결}: VGGT 코드(vggt\_calibrate.py:153--161) 참조.
\end{failbox}


\section{SMPL 피팅 함정}

\begin{failbox}
\textbf{F-4: smplx FK 필수}\\
직접 Forward Kinematics 구현 시 blend shape, joint regressor를 무시하면 MPJPE 1384mm 오차.\\
반드시 \texttt{smplx.create().forward()} 사용.
\end{failbox}

\begin{failbox}
\textbf{F-5: pymomentum import 순서}\\
SAM4Dcap 관련: \texttt{import pymomentum}이 \texttt{import torch}보다 먼저 와야 segfault 방지.
\end{failbox}

\begin{failbox}
\textbf{F-6: cm/m 스케일 불일치}\\
SAM 3D Body는 미터 출력, Meta MHR 변환 코드는 센티미터 가정.\\
\texttt{0.01} 스케일링 팩터를 제거해야 합니다 (SAM4Dcap에서 수정됨).
\end{failbox}

\begin{failbox}
\textbf{F-7: 축-각도 $\pi$ 특이점}\\
투구 가속기에서 어깨 누적 회전 $>$ 180° $\to$ Rodrigues 공식 야코비안 순위 저하.\\
최적화가 진동하거나 발산할 수 있음.\\
\textbf{해결}: 6D 회전 표현 또는 쿼터니언 기반 최적화로 전환 검토.
\end{failbox}


\section{GART 함정}

\begin{failbox}
\textbf{F-8: CUDA 버전 충돌}\\
GART는 CUDA 11.8 필수. 시스템 nvcc 12.4가 우선될 수 있음.\\
\texttt{conda activate gart} 후 \texttt{nvcc --version}이 11.8인지 확인.\\
\textbf{해결}: \texttt{export CUDA\_HOME=\$CONDA\_PREFIX}
\end{failbox}

\begin{failbox}
\textbf{F-9: FLT\_MAX 미정의}\\
\texttt{simple\_knn.cu}에서 \texttt{\#include <cfloat>} 누락 시 컴파일 실패.\\
이미 패치됨 (experiment\_log\_20260414.tex 확인).
\end{failbox}

\begin{failbox}
\textbf{F-10: smplx .J/.A 속성 없음}\\
GART가 사용하는 smplx 버전에서 \texttt{.J}, \texttt{.A} 속성이 없을 수 있음.\\
이미 shim 패치 적용됨 (\texttt{smplx/templates.py}에 수동 J, A 계산 추가).
\end{failbox}

\begin{failbox}
\textbf{F-11: Empty mask crash}\\
프레임에 사람이 없으면 \texttt{get\_bbox}와 \texttt{voxel\_deformer}가 crash.\\
이미 try-except 가드 추가됨.
\end{failbox}


\section{검증 함정}

\begin{failbox}
\textbf{F-12: Vicon 프레임 오프셋}\\
Vicon 데이터는 프레임 512부터 시작 (프레임 0이 아님).\\
RGB는 프레임 0부터 시작.\\
교차상관으로 offset을 찾더라도, Vicon 프레임 번호에 512를 더해야 합니다.
\end{failbox}

\begin{failbox}
\textbf{F-13: 관절 정의 불일치}\\
OpenPose Body25 (25관절) $\neq$ SMPL (24관절) $\neq$ Vicon Plug-in Gait (42마커).\\
비교 시 공통 관절 부분집합 (14개)만 사용해야 합니다.\\
매핑: Neck, RSho, RElb, RWri, LSho, LElb, LWri, RHip, RKne, RAnk, LHip, LKne, LAnk, MidHip.
\end{failbox}


\begin{summary}
\textbf{Claude 실행 체크리스트}:
\begin{enumerate}
\item 서버 접속 전 이 문서의 Chapter 5 (SKILL 1--11)과 Chapter 6 (함정 F-1$\sim$F-13) 숙지
\item 각 SKILL을 순서대로 실행, 검증 기준 통과 확인 후 다음 단계
\item 함정 목록에 해당하는 상황 발생 시 즉시 해결책 참조
\item 품질 게이트 (G1: 7-view 완비, G2: reproj$<$100px, G3: PSNR$>$20) 통과 확인
\item 각 단계 완료 시 결과를 experiment\_log에 기록
\end{enumerate}
\end{summary}

\end{document}
